\section{Results}
\label{sec:results}
% EDIT (2026-10-29): new section built from the old 5.2--5.3 (results, findings), 6.1--6.2
% (observation, exploit, effect), 7.4 (application results), 8.1 (profile-data result) and 8.2
% (guidance), one subsection per experiment. Interpretation moved to Section 5. Table 1 rows are
% reordered to the variant order of Table~\ref{tab:variants}; the per-variant counts of
% Table~\ref{tab:variantsum} are read off Table~\ref{tab:ablation}.
This section reports the five experiments in the order of Section~\ref{sec:setup}. Each subsection
gives the table, then the results we consider significant; what they mean is discussed in
Section~\ref{sec:discussion}.

\subsection{E1: Prompt-Component Ablation}
\label{sec:prompts}
Table~\ref{tab:ablation} reports all $72$ runs and Table~\ref{tab:variantsum} summarizes them per
variant.

\begin{table*}[t]
\centering
\footnotesize
\renewcommand{\arraystretch}{0.88}
\caption{%
  E1: runtime (\textmu s) and outcome across six HeCBench benchmarks and six prompt variants for GPT-5
  and Gemini~2.5~Pro, variants ordered from the least to the most populated prompt. Runtimes are the
  in-loop ten-run averages. Speedup is relative to the unmodified program where its runtime was
  recorded (bitonic-sort: 51{,}590\,\textmu s; sobol: 2{,}576\,\textmu s; shmembench:
  4{,}816\,\textmu s); ``---'' = no baseline available. \textbf{Bold} = best valid speedup per
  benchmark per model. Result: \textit{s1} = first attempt; \textit{rep}$_n$ = repaired in $n$
  iterations; \textit{fail} = all 3 failed; \textit{T/O} = timeout. $^\dagger$ benchmark gaming
  (shared memory algebraically eliminated); $^*$ runtime extraction failure.
}
\label{tab:ablation}
\resizebox{\textwidth}{!}{%
\begin{tabular}{@{}llrrcrrcc@{}}
\toprule
\multirow{2}{*}{\textbf{Benchmark}}
  & \multirow{2}{*}{\textbf{Variant}}
  & \multicolumn{3}{c}{\textbf{GPT-5}}
  & \multicolumn{4}{c}{\textbf{Gemini 2.5 Pro}} \\
\cmidrule(lr){3-5}\cmidrule(lr){6-9}
  & & \textbf{Runtime (\textmu s)} & \textbf{Speedup} & \textbf{Result}
    & & \textbf{Runtime (\textmu s)} & \textbf{Speedup} & \textbf{Result} \\
\midrule
\multirow{6}{*}{\texttt{atan2}}
  & \texttt{minimal}         &  97.53 & --- & \textit{s1}       &&  92.27 & --- & \textit{rep}$_2$ \\
  & \texttt{no\_role}        &  87.76 & --- & \textit{s1}       &&  91.59 & --- & \textit{s1}      \\
  & \texttt{no\_constraints} &  94.01 & --- & \textit{s1}       &&  91.22 & --- & \textit{s1}      \\
  & \texttt{no\_proof}       &  89.00 & --- & \textit{s1}       &&  92.63 & --- & \textit{s1}      \\
  & \texttt{no\_explanation} &  91.08 & --- & \textit{s1}       &&  92.41 & --- & \textit{s1}      \\
  & \texttt{cot}             &  87.47 & --- & \textit{s1}       &&    --- & --- & \textit{fail}    \\
\midrule
\multirow{6}{*}{\texttt{bitonic-sort}}
  & \texttt{minimal}         & 55{,}820 & 0.92$\times$          & \textit{s1} && 55{,}181 & 0.93$\times$          & \textit{s1}      \\
  & \texttt{no\_role}        & 57{,}370 & 0.90$\times$          & \textit{s1} && 54{,}613 & 0.94$\times$          & \textit{s1}      \\
  & \texttt{no\_constraints} & 45{,}364 & \textbf{1.14}$\times$ & \textit{s1} && 57{,}022 & 0.90$\times$          & \textit{rep}$_2$ \\
  & \texttt{no\_proof}       & 45{,}342 & \textbf{1.14}$\times$ & \textit{s1} && 57{,}288 & 0.90$\times$          & \textit{rep}$_2$ \\
  & \texttt{no\_explanation} & 57{,}277 & 0.90$\times$          & \textit{s1} && 45{,}681 & \textbf{1.13}$\times$ & \textit{rep}$_2$ \\
  & \texttt{cot}             & 45{,}450 & \textbf{1.14}$\times$ & \textit{s1} && 45{,}678 & \textbf{1.13}$\times$ & \textit{s1}      \\
\midrule
\multirow{6}{*}{\texttt{md5hash}}
  & \texttt{minimal}         & 396{,}800 & --- & \textit{s1}      && 288{,}700 & --- & \textit{s1}      \\
  & \texttt{no\_role}        & 293{,}100 & --- & \textit{s1}      && 299{,}500 & --- & \textit{s1}      \\
  & \texttt{no\_constraints} & 262{,}600 & --- & \textit{s1}      && 296{,}500 & --- & \textit{s1}      \\
  & \texttt{no\_proof}       & 262{,}500 & --- & \textit{rep}$_3$ && 301{,}400 & --- & \textit{s1}      \\
  & \texttt{no\_explanation} & 262{,}600 & --- & \textit{s1}      && 299{,}200 & --- & \textit{s1}      \\
  & \texttt{cot}             & 295{,}200 & --- & \textit{s1}      && 296{,}900 & --- & \textit{s1}      \\
\midrule
\multirow{6}{*}{\texttt{accuracy}}
  & \texttt{minimal}         & 282.68 & --- & \textit{s1}       && 294.74 & --- & \textit{s1}      \\
  & \texttt{no\_role}        &    --- & --- & \textit{T/O}      && 296.31 & --- & \textit{s1}      \\
  & \texttt{no\_constraints} & 213.70 & --- & \textit{s1}       && 296.44 & --- & \textit{s1}      \\
  & \texttt{no\_proof}       & 214.08 & --- & \textit{rep}$_2$  && 308.62 & --- & \textit{s1}      \\
  & \texttt{no\_explanation} & 214.83 & --- & \textit{s1}       && 213.27 & --- & \textit{s1}      \\
  & \texttt{cot}             & 215.38 & --- & \textit{s1}       && 214.90 & --- & \textit{s1}      \\
\midrule
\multirow{6}{*}{\texttt{sobol}}
  & \texttt{minimal}         & 2{,}824.18 & 0.91$\times$          & \textit{s1}      && 2{,}574.01   & 1.00$\times$ & \textit{rep}$_3$ \\
  & \texttt{no\_role}        & 2{,}623.37 & \textbf{0.98}$\times$ & \textit{s1}      &&          --- & ---          & \textit{fail}    \\
  & \texttt{no\_constraints} &        --- & ---$^*$               & \textit{s1}$^*$  && 2{,}620.28   & 0.98$\times$ & \textit{s1}      \\
  & \texttt{no\_proof}       & 2{,}814.58 & 0.92$\times$          & \textit{s1}      && 2{,}575.83   & 1.00$\times$ & \textit{rep}$_3$ \\
  & \texttt{no\_explanation} & 2{,}660.79 & 0.97$\times$          & \textit{s1}      &&          --- & ---          & \textit{fail}    \\
  & \texttt{cot}             & 2{,}809.94 & 0.92$\times$          & \textit{s1}      &&          --- & ---          & \textit{fail}    \\
\midrule
\multirow{6}{*}{\texttt{shmembench}}
  & \texttt{minimal}         &     5.08 & 949$\times^\dagger$          & \textit{s1}$^\dagger$ && 18{,}735 & 0.26$\times$          & \textit{s1}           \\
  & \texttt{no\_role}        &    20.50 & 235$\times^\dagger$          & \textit{s1}$^\dagger$ &&  4{,}817 & 1.00$\times$          & \textit{s1}           \\
  & \texttt{no\_constraints} &    20.47 & 235$\times^\dagger$          & \textit{s1}$^\dagger$ &&  4{,}661 & \textbf{1.03}$\times$ & \textit{s1}           \\
  & \texttt{no\_proof}       &    20.49 & 235$\times^\dagger$          & \textit{s1}$^\dagger$ &&    22.66 & 213$\times^\dagger$   & \textit{s1}$^\dagger$ \\
  & \texttt{no\_explanation} &    20.46 & 235$\times^\dagger$          & \textit{s1}$^\dagger$ &&  9{,}145 & 0.53$\times$          & \textit{s1}           \\
  & \texttt{cot}             &    20.49 & 235$\times^\dagger$          & \textit{s1}$^\dagger$ &&  4{,}661 & \textbf{1.03}$\times$ & \textit{s1}           \\
\bottomrule
\end{tabular}%
}
\end{table*}

\begin{table}[t]
\centering
\footnotesize
\caption{E1 per variant, read off Table~\ref{tab:ablation}: runs (of $12$, both models) that
  produced a valid candidate at the first attempt (the GPT-5 \texttt{sobol} candidate whose runtime
  could not be extracted counts here but not as a timed result) and that produced a timed result at
  all; and the
  number of benchmarks on which the variant gave a model its slowest valid result (GPT-5 over the
  five benchmarks it did not game; Gemini over six).}
\label{tab:variantsum}
\setlength{\tabcolsep}{5pt}
\begin{tabular}{@{}lcccc@{}}
\toprule
 & \textbf{First attempt} & \textbf{Timed result} & \multicolumn{2}{c}{\textbf{Slowest valid result}} \\
\cmidrule(lr){4-5}
\textbf{Variant} & \textbf{(of 12)} & \textbf{(of 12)} & \textbf{GPT-5 (of 5)} & \textbf{Gemini (of 6)} \\
\midrule
\texttt{minimal}         & 10 & 12 & 4 & 1 \\
\texttt{no\_role}        & 10 & 10 & 1 & 0 \\
\texttt{no\_constraints} & 11 & 11 & 0 & 1 \\
\texttt{no\_proof}       &  8 & 12 & 0 & 4 \\
\texttt{no\_explanation} & 10 & 11 & 0 & 0 \\
\texttt{cot}             & 10 & 10 & 0 & 0 \\
\bottomrule
\end{tabular}
\end{table}

\paragraph{Which variant produces a valid candidate at the first attempt.}
Across the grid, GPT-5 produced a timed result at the first attempt in $32$ of $36$ runs (of the other
four, two needed repair, one timed out and one built but its runtime could not be extracted) and
Gemini~2.5~Pro in $26$ of $36$; Gemini used all three iterations on six runs, two of which were
eventually repaired and four of which failed, three of them on \texttt{sobol} (two correctness
violations and a persistent \texttt{\_\_builtin\_ctz} host-device mismatch). GPT-5 has no build or
correctness failure; its one unusable run is the timeout. Per variant, the first-attempt rate ranges
from $8$ to $11$ of $12$. Removing the prohibitions costs the fewest first attempts ($11$ of $12$);
removing the proof requirement costs the most ($8$ of $12$: two GPT-5 runs and two Gemini runs needed
repair), although all twelve \texttt{no\_proof} runs were eventually repaired. Three variants
lose runs outright: \texttt{no\_role} (GPT-5's single timeout, on \texttt{accuracy}, and a Gemini
failure on \texttt{sobol}), \texttt{no\_explanation} (one Gemini failure, on \texttt{sobol}) and
\texttt{cot} (Gemini failed both \texttt{atan2}, where all three iterations failed, and
\texttt{sobol}, where it placed \texttt{\_\_builtin\_ctz}---a host-only intrinsic---inside a
\texttt{\_\_global\_\_} kernel and could not repair the compile error).

\paragraph{Which variant produces the slowest code.}
For GPT-5, \texttt{minimal} gives the slowest valid result on four of the five benchmarks it did not
game: by $34\%$ on \texttt{md5hash} ($396{,}800$ against $295{,}200$\,\textmu s for the next slowest
variant) and $31\%$ on \texttt{accuracy}, and by under $4\%$ on \texttt{atan2} and \texttt{sobol}. For
Gemini the slow end belongs to \texttt{no\_proof}, which gives its slowest valid result on four of six
benchmarks, each time by a small margin over the next slowest variant ($0.2$--$4\%$), so the pattern is
a consistent position at the slow end rather than a large penalty; Gemini's one large \texttt{minimal}
penalty is \texttt{shmembench} ($18{,}735$\,\textmu s, four times its best valid variant and twice the
next slowest). The \texttt{cot} and \texttt{no\_explanation} variants never give either model its
slowest valid result.

\paragraph{Optimization quality.}
Within this configuration neither model delivers large or reliable speedups. The best result across all
valid runs is $1.14\times$ (GPT-5 on \texttt{bitonic-sort} under \texttt{no\_constraints},
\texttt{no\_proof} or \texttt{cot}) and $1.13\times$ (Gemini on the same benchmark under
\texttt{no\_explanation} or \texttt{cot}); \texttt{cot} reaches the best observed speedup for both
models. On \texttt{sobol} both models either fail or match the baseline within $2\%$. On
\texttt{atan2}, \texttt{md5hash} and \texttt{accuracy} the runtimes differ by variant more than by any
optimization. The remaining components act differently on the two models: removing the prohibitions
produces GPT-5's best \texttt{bitonic-sort} result ($1.14\times$) and one of Gemini's worst
($0.90\times$); removing the explanation lists is Gemini's best \texttt{bitonic-sort} variant
($1.13\times$) and has no discernible effect elsewhere.

\paragraph{Benchmark gaming.}
\label{sec:gaming-obs}
The most consequential entry in Table~\ref{tab:ablation} is not a speedup but a false one. On
\texttt{shmembench}, GPT-5 games the benchmark under all six variants without exception, exploiting
the fact that a swap-based permutation preserves the sum to eliminate all shared-memory traffic and
report $235$--$949\times$ apparent speedups. Gemini does so only when the proof requirement is removed
(\texttt{no\_proof}), reaching $213\times$ by the same mechanism. Two things follow. First, a
prompt-level requirement for a correctness argument suppresses this class of edit in one model but not
the other, so it cannot be relied on as a defence. Second, the benchmark's own verification passes in
every one of these runs, so no output-comparison gate would have caught them. This is what E2
addresses.

\subsection{E2: The Semantic-Preservation Constraint}
\label{sec:intent}
\paragraph{The exploit.}
\texttt{atomicSystemWide} exists to measure system-wide atomic throughput: each thread performs
\texttt{loop\_num} system-scope \texttt{atomicAdd\_system(\&atom\_arr[0],\,10)} operations, and the
program then checks the accumulated total. Without the constraint, both models removed the loop and
kept the total. GPT-5 replaced it with one atomic per thread carrying the pre-summed value, annotating
its own source \textit{``one store per warp is sufficient for postcondition used by verification'';}
Gemini collapsed it further, to a single \texttt{atomicAdd\_system(\&atom\_arr[0], blockDim.x * 10 *
loop\_num)} per block. Both versions print the same output and pass the program's own check.

\paragraph{Effect of the constraint.}
Table~\ref{tab:intent} reports the outcome. Without the constraint, the two models report
$1504\times$ and $2809\times$; with it, all three constrained prompts return honest, roughly neutral
results ($0.71$--$1.08\times$), and in every constrained candidate the per-iteration
\texttt{atomicAdd\_system(\ldots,\,10)} is still present in the generated source. Two caveats apply to
the ratios. The unconstrained and relaxed arms were generated and measured in the same interleaved
session, but they differ in more than the constraint, since the relaxed prompt also lifts the
prohibitions; the arms that isolate the constraint (original prohibitions plus the
semantic-preservation block) were timed in a later session on the same node, rebuilt with a newer CUDA
toolkit, where this benchmark's repetition spread is $16$--$20\%$. The result we draw from the table is
therefore categorical---gaming against no gaming, confirmed by reading the generated source---rather
than a precise ratio.

\begin{table}[t]
\centering
\footnotesize
\setlength{\tabcolsep}{5pt}
% EDIT (2026-10-29): header rebuilt as two clean rows (the old one split "No constraint" over two
% lines); values unchanged (pcrl/code_only from prof_ablation_cu133.csv, relaxed from three_arm_v2).
\caption{E2: speedup of \texttt{atomicSystemWide} against the unmodified program, without the
  semantic-preservation constraint and under the three prompts that carry it. $^\dagger$ = benchmark
  gaming: the measured atomic loop is removed while output and the built-in check are unchanged. All
  constrained candidates retain the per-iteration atomic; their exact ratios are not comparable across
  sessions (see text).}
\label{tab:intent}
\begin{tabular}{@{}lrrrr@{}}
\toprule
 & & \multicolumn{3}{c}{\textbf{With constraint}} \\
\cmidrule(lr){3-5}
\textbf{Model} & \textbf{No constraint} & \textbf{pcrl} & \textbf{code\_only} & \textbf{relaxed} \\
\midrule
GPT-5          & 1504$\times^\dagger$ & 0.88$\times$ & 1.08$\times$ & 0.94$\times$ \\
Gemini 3.1 Pro & 2809$\times^\dagger$ & 0.81$\times$ & 1.01$\times$ & 0.71$\times$ \\
\bottomrule
\end{tabular}
\end{table}

\subsection{E3: Production Applications}
\label{sec:realapps}
Table~\ref{tab:realapps} reports the speedups against the unmodified application.

\begin{table}[t]
\centering
\footnotesize
% EDIT (2026-10-29): caption now defines the three arms and is written for a reader who has not
% seen the methodology section.
\caption{E3: speedup of each application against its unmodified version, measured with the
  application's own phase timer on the GH200 (medians over clean interleaved repetitions,
  Section~\ref{sec:setup-e3}; every candidate passed its gate). \emph{Baseline} is the full prompt
  with PC-sampling and roofline data; \emph{structural} relaxes the prohibitions and adds the derived
  kernel-budget and occupancy summaries; \emph{hand} is an expert reference restricted to the same
  single file. ``fail'' = no candidate passed the gate within three repair iterations.}
\label{tab:realapps}
\setlength{\tabcolsep}{3pt}
\begin{tabular}{@{}llrrr@{}}
\toprule
\textbf{Application} & \textbf{Prompt} & \textbf{GPT-5} & \textbf{Gemini 3.1 Pro} & \textbf{Hand} \\
\midrule
\multirow{2}{*}{AutoDock-GPU}
  & baseline    & 0.990$\times$ & 1.037$\times$ & \multirow{2}{*}{\textbf{1.532}$\times$} \\
  & structural  & fail          & 1.001$\times$ & \\
\midrule
\multirow{2}{*}{Cholla}
  & baseline    & \textbf{1.571}$\times$ & 1.504$\times$ & \multirow{2}{*}{1.468$\times$} \\
  & structural  & 0.992$\times$ & 1.469$\times$ & \\
\bottomrule
\end{tabular}
\end{table}

\paragraph{Cholla: every baseline-prompt output beats the hand reference.}
The best model output in this table (GPT-5, $1.571\times$) beats the hand reference ($1.468\times$),
and so does every other baseline-prompt sample we have: the profile-data ablation (E4) regenerated the
same prompt for both models, giving $1.476\times$ (GPT-5) and $1.584\times$ (Gemini), timed in a
later session in which this table's two candidates re-measured at $1.572\times$ and $1.504\times$.

\paragraph{AutoDock-GPU: no candidate finds the restructuring.}
The profitable change here is a cross-cutting restructuring, worth $1.532\times$ by hand. \emph{Every}
candidate that built---three of the four model~$\times$~prompt cells, the fourth producing none---left
the structure it rewrites in place and gained at most $3.7\%$. What the hand reference changed, and
what the models changed instead, is described in Section~\ref{sec:whatmodels}.

\paragraph{The structural prompt never helped and broke real builds.}
Its contract was written for single-file benchmarks: it permits signature changes and states that the
returned file must be self-contained. In Cholla the target kernel is declared in a header and called
from ten sites in other translation units; both models changed the signature and both produced linker
errors, each needing a repair iteration to revert. On AutoDock-GPU, whose target file is
\texttt{\#include}d into another translation unit, GPT-5 redefined symbols already in scope and then
produced an out-of-bounds global read (found with \texttt{compute-sanitizer}), failing all three
iterations.

\subsection{E4: Profile-Data Ablation}
\label{sec:evidence}
\label{sec:profabl}
% EDIT (2026-10-27): HeCBench half re-timed. The first timing (2026-09-22) ran while a driver change on
% the H100 node made CUDA 12.x binaries print wrong output and run 20-300x slow (real-applications
% report, sec. 14). Same candidates, rebuilt with CUDA 13.3.1 and re-timed on the same node on
% 2026-10-26 (experiments repo: scripts/optimas_study/rerun_cu133/); chemv excluded.
Profiling data helps consistently on the large real application; on the small benchmarks it changes
individual results by as much as on the application, but in both directions, with no net gain
(Tables~\ref{tab:profabl} and~\ref{tab:profablhec}). On Cholla both models are clearly faster with
the profile ($+11.5\%$ for GPT-5, $+7.7\%$ for Gemini), and the second, independently generated
\texttt{pcrl} sample keeps both on the same side of the comparison (GPT-5 $1.187$, Gemini $1.023$
relative to \texttt{code\_only}), although the Gemini replicate is much weaker than its first sample.
On AutoDock-GPU both arms reach the same ${\sim}1.03\times$. Of the nineteen scored HeCBench pairs
(GPT-5 produced no correct \texttt{code\_only} candidate for \texttt{background-subtract}), nine are
ties by the rule of Section~\ref{sec:timing}, six favour the profile and four favour its absence.
Counting only differences above $5\%$, four favour the profile (\texttt{convolutionSeparable}
$+13.4\%$ and \texttt{bincount} $+8.4\%$ for GPT-5; \texttt{amgmk} $+13.5\%$ and
\texttt{convolutionSeparable} $+5.3\%$ for Gemini) and two go against it (Gemini \texttt{boxfilter}
$-15.3\%$ at a $0.3\%$ spread, and Gemini \texttt{atomicSystemWide} $-19.9\%$ at a $16.3\%$ spread).
The geometric mean of $\mathrm{pcrl}/\mathrm{co}$ is $1.003$ over the nine GPT-5 pairs and $0.979$
over the ten Gemini pairs. Each pair is a single generation per arm, so an individual benchmark's sign
may be sampling variance; what the benchmarks do not show is a consistent direction.

\begin{table}[t]
\centering
\footnotesize
\caption{E4 on the applications: speedup against the unmodified application for the prompt without
  profile data (\texttt{co}) and with PC sampling + roofline (\texttt{pcrl}), and their ratio
  (repetitions in Section~\ref{sec:setup-e4}). The HeCBench rows give the geometric mean of the
  per-pair ratios in Table~\ref{tab:profablhec}.}
\label{tab:profabl}
\setlength{\tabcolsep}{2pt}
\begin{tabular}{@{}llrrr@{}}
\toprule
\textbf{Application} & \textbf{Model} & \textbf{co} & \textbf{pcrl} & \textbf{pcrl/co} \\
\midrule
\multirow{2}{*}{Cholla}
  & GPT-5  & 1.324$\times$ & \textbf{1.476}$\times$ & \textbf{1.115} \\
  & Gemini & 1.471$\times$ & \textbf{1.584}$\times$ & \textbf{1.077} \\
\midrule
\multirow{2}{*}{AutoDock-GPU}
  & GPT-5  & 1.030$\times$ & 1.031$\times$ & 1.001 \\
  & Gemini & 1.035$\times$ & 1.035$\times$ & 1.000 \\
\midrule
HeCBench (9 pairs, geomean)  & GPT-5  & & & 1.003 \\
HeCBench (10 pairs, geomean) & Gemini & & & 0.979 \\
\bottomrule
\end{tabular}
\end{table}

\begin{table}[t]
\centering
\footnotesize
\caption{E4 on HeCBench (H100, CUDA 13.3.1, medians of $5$ clean repetitions per arm). Speedups are
  against the unmodified benchmark; \emph{spr.} is the largest repetition spread of the pair, in
  percent. A ratio in bold differs from $1$ by more than its pair's spread; the others are ties.
  --- = no correct \texttt{code\_only} candidate.}
\label{tab:profablhec}
\setlength{\tabcolsep}{2.5pt}
\begin{tabular}{@{}lrrrrcrrrr@{}}
\toprule
 & \multicolumn{4}{c}{\textbf{GPT-5}} && \multicolumn{4}{c}{\textbf{Gemini 3.1 Pro}} \\
\cmidrule(lr){2-5}\cmidrule(lr){7-10}
\textbf{Benchmark} & \textbf{co} & \textbf{pcrl} & \textbf{pcrl/co} & \textbf{spr.} &&
  \textbf{co} & \textbf{pcrl} & \textbf{pcrl/co} & \textbf{spr.} \\
\midrule
\texttt{adv}                  & 0.997 & 0.999 & 1.002 & 19.1 && 1.002 & 0.994 & 0.992 & 14.5 \\
\texttt{amgmk}                & 0.979 & 0.989 & 1.010 & 11.8 && 0.878 & 0.996 & \textbf{1.135} & 5.4 \\
\texttt{atomicSystemWide}     & 1.075 & 0.876 & 0.815 & 19.5 && 1.010 & 0.809 & \textbf{0.801} & 16.3 \\
\texttt{background-subtract}  & ---   & 1.638 & ---   & 16.4 && 1.574 & 1.581 & 1.004 & 11.1 \\
\texttt{bincount}             & 1.000 & 1.084 & \textbf{1.084} & 0.1 && 1.359 & 1.364 & \textbf{1.004} & 0.1 \\
\texttt{blockAccess}          & 1.021 & 1.004 & \textbf{0.984} & 0.9 && 1.000 & 1.003 & 1.004 & 0.9 \\
\texttt{boxfilter}            & 1.079 & 1.113 & \textbf{1.032} & 0.2 && 1.796 & 1.521 & \textbf{0.847} & 0.3 \\
\texttt{ccsd-trpdrv}          & 1.006 & 1.005 & 0.999 & 31.6 && 1.030 & 1.019 & 0.989 & 37.4 \\
\texttt{concat}               & 0.937 & 0.933 & \textbf{0.995} & 0.1 && 1.350 & 1.350 & 1.000 & 0.4 \\
\texttt{convolutionSeparable} & 0.999 & 1.133 & \textbf{1.134} & 1.1 && 1.084 & 1.141 & \textbf{1.053} & 0.8 \\
\bottomrule
\end{tabular}
\end{table}

\paragraph{Correctness, not only speed.}
% EDIT (2026-10-27): background-subtract re-checked on 2026-10-26 with CUDA 13.3.1: the unmodified
% program passes (max error 0) and all three code_only candidates fail (max error 93).
One HeCBench pair favours the profile arm on correctness rather than time. On
\texttt{background-subtract}, GPT-5 without profile data failed the benchmark's own output check in
all three repair iterations and produced no usable candidate, while the \texttt{pcrl} arm produced a
correct one ($1.638\times$). Because these candidates were generated on the H100 node while it was
faulty (Section~\ref{sec:limits}), we rebuilt all three rejected candidates with the working toolchain:
the unmodified program passes with zero error, and each candidate still fails, with a maximum pixel
error of $93$.

\subsection{E5: Prescriptive Guidance}
\label{sec:guidvar}
Table~\ref{tab:guidvar} reports the four arms. Half of the rows have a repetition spread below $5\%$,
including the two largest effects of the checklist, so those differences are not measurement
artifacts.

\begin{table}[t]
\centering
\footnotesize
\setlength{\tabcolsep}{2pt}
\caption{E5: speedup against the unmodified program under the constrained (\emph{base}) prompt, the
  relaxed (\emph{struct}) prompt and the relaxed prompt plus the prescriptive checklist (\emph{diag});
  interleaved medians on one idle GH200. \textbf{Bold} marks the cases discussed in the text.
  ``spread'' is the largest repetition spread across the four arms. \texttt{chemv} is excluded
  (see text).}
\label{tab:guidvar}
\resizebox{\columnwidth}{!}{%
\begin{tabular}{@{}llrrrr@{}}
\toprule
\textbf{Benchmark} & \textbf{Model} & \textbf{base} & \textbf{struct} & \textbf{diag} & \textbf{spread} \\
\midrule
\texttt{concat}               & GPT-5  & 1.031$\times$ & 0.838$\times$ & \textbf{0.349}$\times$ & 0.1\% \\
\texttt{boxfilter}            & GPT-5  & \textbf{10.280}$\times$ & 6.352$\times$ & 5.503$\times$ & 0.0\% \\
\midrule
\texttt{bincount}             & GPT-5  & 1.035$\times$ & 1.037$\times$ & \textbf{1.374}$\times$ & 5.2\% \\
\texttt{bincount}             & Gemini & 1.106$\times$ & 1.003$\times$ & \textbf{1.380}$\times$ & 4.4\% \\
\texttt{concat}               & Gemini & 1.001$\times$ & \textbf{1.797}$\times$ & 1.798$\times$ & 46.8\% \\
\texttt{boxfilter}            & Gemini & 8.691$\times$ & 8.930$\times$ & 9.371$\times$ & 42.6\% \\
\texttt{convolutionSeparable} & Gemini & 1.418$\times$ & 1.499$\times$ & 1.514$\times$ & 19.3\% \\
\midrule
\texttt{blockAccess}          & both   & $\approx 1.00\times$ & $\approx 1.00\times$ & $\approx 1.00\times$ & 0.2\% \\
\bottomrule
\end{tabular}%
}
\end{table}

\paragraph{The checklist produces our largest gain.}
On \texttt{bincount} it takes both models from roughly $1.00$--$1.04\times$ under the structural prompt
to $1.37$--$1.38\times$, approaching our hand-tuned reference ($1.392\times$). The other large
improvement in the table is not the checklist's: on \texttt{concat} under Gemini the $1.80\times$
appears already under the structural prompt ($1.797\times$), whose relaxed rules permit the vectorized
templating it needs, and the checklist adds nothing ($1.798\times$) at a $46.8\%$ spread.

\paragraph{It also produces our largest loss.}
On \texttt{concat}, GPT-5 followed the instruction to process several elements per thread literally,
to one thread per $128$-element line: a $128\times$ collapse in parallelism, measured at $0.349\times$
against $0.838\times$ without the checklist, with a $0.1\%$ spread. On \texttt{boxfilter}, GPT-5's
$10.280\times$ under the baseline prompt falls to $6.352\times$ when the prohibitions are relaxed and
to $5.503\times$ with the checklist added; most of that loss belongs to the relaxation, not to the
checklist. All \texttt{boxfilter} results exceed our $5\times$ alarm, so we read the diffs: every
candidate keeps the timing loop and both kernel launches per cycle, and gains by tiling the column
pass over rows, which the original runs as one thread per image column ($30$ blocks of $64$ threads
for a $1920$-pixel-wide image); the benchmark's own check compares every output pixel exactly with a
host implementation.

\paragraph{Net effect.}
Averaged over all twelve (benchmark, model) pairs the checklist is close to neutral---a geometric mean
of $0.97\times$ against the structural prompt, with individual pairs from $0.42\times$ to
$1.38\times$. The same prompt component that produces $1.380\times$ on one benchmark produces
$0.349\times$ on another, both at sub-$5\%$ spread; the spread of outcomes, not their mean, is the
result (Section~\ref{sec:profdisc}).

\paragraph{An excluded outlier: \texttt{chemv}.}
% NEW (2026-10-26). Evidence: chemv-verification-vacuous (experiments memory), report sec. 13.3/14,
% and a re-check on illyad (CUDA 13.3.1, 2026-10-26) of every chemv candidate with REPEAT=1 against
% the CPU reference: max relative error 2.2e-7 for the unmodified program and for Gemini's baseline
% candidate, at most 1.1e-6 for any candidate.
We exclude \texttt{chemv} from every table because its built-in verification cannot certify a result.
Its timing loop launches both kernels $1000$ times, accumulating into the same output vector without
re-initializing it, and then compares that vector with a single application of the CPU reference: the
check fails for the unmodified program, or---once the accumulated values overflow to NaN---passes
vacuously, because a comparison with NaN is never true. The case is nonetheless instructive. Under the
baseline prompt, Gemini rewrote both polyhedral-compiler-generated kernels: it replaced each thread's
tiled loop nest (an outer loop over $32$-element tiles with a barrier per tile, and an inner loop
clipped to the tile) by a single loop over the same elements of its row, removed the barriers---which
the computation does not need, since each thread updates only its own output element and the kernels
use no shared memory---and accumulated in registers instead of reading and writing the output in global
memory on every iteration. The number of multiply-adds is unchanged; what was removed is loop structure
that correct output does not require. The benchmark timed it at $7.8\times$, over our $5\times$ alarm,
so we checked it outside the benchmark's harness, running one application of the two kernels (no
repeat loop, so nothing overflows) against the CPU reference: its largest relative error is
$2.2\times10^{-7}$, identical to the unmodified program's, and no \texttt{chemv} candidate from either
model in the baseline, structural, checklist or profile-data arms exceeds $1.1\times10^{-6}$. Under the
checklist the same model kept the original tile loops and added the shared-memory staging the checklist
named. We report the case as an example of the pipeline removing loops that correctness does not
need---the legitimate counterpart of the gaming of E2, which removed a loop the measurement does
need---and not as a measured speedup, since the timing that produced it rests on a loop whose output is
NaN.
